\section{Strong-Weak Involution}
\label{sec:strong_weak_involution}

\paragraph{Involution on the crossover coordinate.}
Let
\begin{equation}
  y = \log \chi, \qquad \chi > 0.
\end{equation}
Define
\begin{equation}
  \mathcal{D}_{\mathrm{sw}}: y \mapsto -y
  \quad \Longleftrightarrow \quad
  \chi \mapsto \chi^{-1}.
\end{equation}

\begin{theorem}[Strong-weak involution]
\label{thm:strong_weak_involution}
Assume one-channel scaling
\begin{equation}
  \chi(\beta,g,\vartheta)=g^2\chi_0(\beta,\vartheta), \qquad \chi_0(\beta,\vartheta)>0.
\end{equation}
Define the self-dual scale
\begin{equation}
  g_\star(\beta,\vartheta):=\chi_0(\beta,\vartheta)^{-1/2},
\end{equation}
and the dual representative
\begin{equation}
  g^\vee(\beta,\vartheta;g):=\frac{g_\star(\beta,\vartheta)^2}{g}
  =\frac{1}{\chi_0(\beta,\vartheta)\,g}.
\end{equation}
Then:
\begin{align}
  \chi(\beta,g^\vee,\vartheta) &= \chi(\beta,g,\vartheta)^{-1}, \\
  (g^\vee)^\vee &= g, \\
  y(\beta,g^\vee,\vartheta) &= -y(\beta,g,\vartheta).
\end{align}
\end{theorem}
\begin{proof}
By substitution,
\begin{equation}
  \chi(\beta,g^\vee,\vartheta)
  =\chi_0\left(\frac{1}{\chi_0 g}\right)^2
  =\frac{1}{\chi_0 g^2}
  =\chi(\beta,g,\vartheta)^{-1}.
\end{equation}
Applying the same transformation twice gives
\begin{equation}
  (g^\vee)^\vee=\frac{1}{\chi_0 g^\vee}
  =\frac{1}{\chi_0}\,\chi_0 g
  =g.
\end{equation}
Taking logarithms yields \(y\mapsto -y\).
\end{proof}

\begin{proposition}[Asymptotic exchange]
\label{prop:asymptotic_exchange}
Let \(Q\) be admissible in the sense of Definition~\ref{def:admissible_class}, and suppose the
\(\chi\)-dependent building blocks \(\phi_j\) have weak/strong asymptotic expansions at
\(\chi\to 0\) and \(\chi\to\infty\). Then the strong-\(\chi\) asymptotic expansion of \(Q\) at
\((\beta,g,\vartheta)\) is obtained by evaluating the weak-\(\chi\) expansion at the dual
representative \((\beta,g^\vee,\vartheta)\), and conversely.
\end{proposition}
\begin{proof}
By Theorem~\ref{thm:strong_weak_involution},
\(\chi(\beta,g^\vee,\vartheta)=\chi(\beta,g,\vartheta)^{-1}\). Therefore any weak expansion
\(\phi_j(\chi)\sim \sum_n a_{j,n}\chi^n\) at \(g^\vee\) becomes
\(\sum_n a_{j,n}\chi(\beta,g,\vartheta)^{-n}\), which is a strong expansion at \(g\). The reverse
direction is identical.
\end{proof}

\paragraph{Interpretation boundary.}
The proposition transfers asymptotic control between representatives. It is not a claim that
\(Q(\beta,g,\vartheta)=Q(\beta,g^\vee,\vartheta)\) pointwise.
